\section{Uniform control at logarithmic zero}
\label{sec:near-zero-trace}

The causal and anticausal test classes do not give a smooth
partition of an arbitrary test through zero.  The following
argument controls the omitted neighborhood in the actual
trace norm.  It uses the proved weighted spectral synthesis,
but needs no uniform minimum-modulus estimate in a moving norm.

\paragraph{MA-NZ-001 (a uniform actual trace-ideal estimate near zero).}
Use the rational family of MA-GW-001--003, or either character
family of MA-GW-004.  Fix $0<\alpha<1/2$, $s\ge0$ and the
single physical space
\[
 H=L^2(\mathbb R,w(x)\,dx),\qquad
 w(x)=e^{2\alpha|x|}(1+x^2)^{s/2}.
 \tag{MA-NZ.1}
\]
Let $c\le0$, let $J_c\subset H_c=\mathbf1_{[c,\infty)}H$
be the actual supported PW range, and let $P_c^Z$ project
orthogonally onto $Z_c=H_c\ominus J_c$, extended by zero.
For $g\in C_c^\infty(\mathbb R)$ with $g(0)=0$, choose
$\chi\in C_c^\infty((-2,2))$ equal to one on $[-1,1]$, and set
\[
                 r_\varepsilon(t)=g(t)\chi(t/\varepsilon).
 \tag{MA-NZ.2}
\]
Then
\[
       \lim_{\varepsilon\downarrow0}
       \sup_{c\le0}\|C_{r_\varepsilon}P_c^Z\|_1=0.
 \tag{MA-NZ.3}
\]
All the operators in this assertion act on the actual $H$.
The result holds at every polynomial weight, with the exact
upper-edge truncations of MA-GW.

\begin{proof}
\emph{1. Uniform dissipativity and an actual Hilbert--Schmidt resolvent.}
On the physical supported half-line use the positive translations
$U_tf(x)=\mathbf1_{x\ge c+t}f(x-t)$, $t\ge0$.
The logarithm of $w$ is globally Lipschitz and satisfies
\[
 (\log w)'(x)\le2\alpha+s/2\quad\text{almost everywhere}.
\]
Consequently $\|U_t\|\le e^{\omega t}$ for
$\omega=\alpha+s/4$, independently of $c$.
The actual quotient semigroup has the same bound.
Its Hilbert adjoint on $Z_c$ also has the same bound.
The full half-line adjoint generator is
\[
 \begin{aligned}
 \mathcal L_cv&=w^{-1}(wv)',\\
 \mathcal D(\mathcal L_c^2)
   &=\{v:(wv)^{(j)}\in L^2((c,\infty),w^{-1}dx),\ 0\le j\le2\}.
 \end{aligned}
 \tag{MA-NZ.4}
\]
There is no boundary condition at $c$.
This is the graph domain proved in MA-GW.37--.38 or
MA-GW.78, expressed in physical coordinates.
Let $B_c$ be its generator restricted to the adjoint-invariant
space $Z_c$.  Choose a fixed real $\nu>\omega$, put
$a_0=\nu-\omega>0$, and extend
\[
                   R_c=(\nu-B_c)^{-1}P_c^Z
 \tag{MA-NZ.5}
\]
by zero on $Z_c^\perp$.
The actual vectorwise Laplace resolvent gives
$\|R_c\|\le a_0^{-1}$.

Differentiating the quasicontractive semigroup norm at zero
on generator-domain vectors gives
$\operatorname{Re}\langle B_cv,v\rangle\le\omega\|v\|^2$.
The same holds for the adjoint generator, since its semigroup
has the same norm bound.  Applying these two inequalities
to the two resolvents proves, on $H$,
\[
 \begin{gathered}
 0\le a_0R_c^*R_c\le\operatorname{Re}R_c,
 \qquad 0\le a_0R_cR_c^*\le\operatorname{Re}R_c,\\
 \operatorname{Re}R_c=(R_c+R_c^*)/2.
 \end{gathered}
 \tag{MA-NZ.6}
\]
For example, with $v=R_cf$ in $Z_c$,
$\operatorname{Re}\langle R_cf,f\rangle
 =\nu\|v\|^2-\operatorname{Re}\langle v,B_cv\rangle
 \ge a_0\|R_cf\|^2$.

Write $\lambda_\rho=\rho-1/2$ and let $r_\rho^+$ be the
actual retained half-line lengths from MA-GW, in the indicated
character sector.  The exact kernel assertion there implies
that $Z_c$ is the closed span of the Riesz representatives
\[
  \mathbf1_{[c,\infty)}(x)\,w(x)^{-1}x^j
             e^{\overline{\lambda_\rho}x},
                     \qquad 0\le j<r_\rho^+.
 \tag{MA-NZ.7}
\]
This proves the needed completeness of the adjoint blocks;
it is not inferred merely from completeness of an unrelated
primal spectral family.
Let $F_{T,c}$ be the span of these vectors with
$|\operatorname{Im}\rho|\le T$, and let $E_{T,c}$ be its
orthogonal projection in $H$.
These are finite-dimensional invariant spaces for $B_c$
and $R_c$, have dimensions independent of $c$, and increase
densely to $Z_c$.
Their eigenvalues for $B_c$ are
$\overline{\lambda_\rho}$ with the stated lengths.

Since the trace on a finite invariant space counts those
eigenvalues, even with nontrivial Jordan blocks,
\[
 \begin{aligned}
 \operatorname{Tr}(E_{T,c}\operatorname{Re}R_cE_{T,c})
   &=\sum_{|\operatorname{Im}\rho|\le T}
       r_\rho^+\operatorname{Re}(\nu-\lambda_\rho)^{-1},\\
 S_\nu&=\sum_\rho
       r_\rho^+\operatorname{Re}(\nu-\lambda_\rho)^{-1}<\infty.
 \end{aligned}
 \tag{MA-NZ.8}
\]
Every term is nonnegative.  The containing strip bounds
each resolvent factor by a constant times
$(1+|\operatorname{Im}\rho|^2)^{-1}$, and the
multiplicity-weighted $O(\log T)$ local count proves
summability.
Choose orthonormal bases successively along these nested
finite spaces.  Positivity and (MA-NZ.8) show that
$(\operatorname{Re}R_c)^{1/2}$ is Hilbert--Schmidt;
hence $\operatorname{Re}R_c$ is trace class with trace
$S_\nu$, independent of $c$.
Equation (MA-NZ.6) now gives $R_c\in\mathcal S_2$ and
$\|R_c\|_2^2\le S_\nu/a_0$.

Let $S_\nu(T)$ be the tail of the sum in (MA-NZ.8).
Compressing both positive inequalities in (MA-NZ.6)
to $1-E_{T,c}$ gives both uniform tails
\[
 \begin{aligned}
 \|R_c(1-E_{T,c})\|_2^2&\le S_\nu(T)/a_0,\\
 \|(1-E_{T,c})R_c\|_2^2&\le S_\nu(T)/a_0,\\
 \|R_c^2-R_cE_{T,c}R_c\|_1
       &\le a_0^{-1}\sqrt{S_\nu S_\nu(T)}.
 \end{aligned}
 \tag{MA-NZ.9}
\]
In particular $R_c^2$ is trace class with a uniform
finite-block approximation.  The full half-line resolvent
preserves $Z_c$, so
$R_c^2=(\nu-\mathcal L_c)^{-2}P_c^Z$ as an actual operator.
The identities
$\mathcal L_cR_c^2=\nu R_c^2-R_c$ and
$\mathcal L_c^2R_c^2=\nu^2R_c^2-2\nu R_c+P_c^Z$
also give a uniform bound into the second graph domain.

\emph{2. Finite-block regularity without a moving-norm estimate.}
For each fixed $T$, the norm of $B_c|_{F_{T,c}}$ is
bounded independently of $c$.
Here is an elementary finite-dimensional proof.
A complex matrix is unitarily upper triangularizable:
choose a unit eigenvector, decompose its orthogonal
complement, and apply induction to the compression on
that complement.  The resulting triangular matrix has
the eigenvalues on its diagonal.
In this basis quasidissipativity says that
$\omega I-\operatorname{Re}B_c$ is positive semidefinite.
Its two-by-two principal minors therefore give, for $i<j$,
\[
 |(B_c)_{ij}|
   \le2\sqrt{(\omega-\operatorname{Re}\lambda_i)
                     (\omega-\operatorname{Re}\lambda_j)}.
 \tag{MA-NZ.10}
\]
The real parts lie in the fixed containing strip, the
heights are bounded by $T$, and the dimension is fixed.
Bounding the norm by the diagonal norm plus the
Hilbert--Schmidt norm of the strict upper triangle
gives the claimed bound $C_T$.
This uses neither a basis condition number nor zero separation.

Put $\ell=w'/w$ on the two open half-lines.  It is bounded
by $2\alpha+s/2$ in absolute value.
For a vector $v\in F_{T,c}$, the identity
$v'=B_cv-\ell v$ bounds its weighted first derivative
by $C_T\|v\|$ on $(c,\infty)$.
The one-dimensional Sobolev trace inequality on $[c,c+1]$
also gives
\[
             \sqrt{w(c)}\,|v(c+)|\le C_T\|v\|_H,
 \tag{MA-NZ.11}
\]
uniformly in $c$, because the ratio $w(c+t)/w(c)$
is uniformly bounded above and below for $0\le t\le1$.
Extend $v$ by zero to the full line.
The fundamental theorem of calculus bounds the translated
difference away from the endpoint by $C_T|t|\|v\|$.
On the endpoint interval of length $|t|$, (MA-NZ.11)
and the same derivative bound give $C_T\sqrt{|t|}\|v\|$.
For negative translations apply the bounded inverse shift
to the positive-translation estimate.
Thus, for $|t|\le1$,
\[
      \|U_tv-v\|_H\le C_T\sqrt{|t|}\,\|v\|_H
                      \qquad(v\in F_{T,c},\ c\le0).
 \tag{MA-NZ.12}
\]
Here $U_t$ is the full-line translation; its small-shift
norm is uniformly bounded by the same Lipschitz weight bound.

\emph{3. The actual second-order ambient factorization.}
For any smooth compact $r$ write
$a=w^{-1}$ and $b=wv$ on $(c,\infty)$.
For $v\in\mathcal D(\mathcal L_c^2)$, two integrations
by parts in
$\int_c^\infty r(x-y)a(y)(\nu-\partial_y)^2b(y)dy$
give a bounded bulk operator
\[
 \mathcal B_r
   =C_{r''}-2C_{r'}M_{\nu-\ell}
          +C_rM_{(\nu-\ell)^2-\ell'_{\rm reg}}.
 \tag{MA-NZ.13}
\]
The bounded regular derivative is
$\ell'_{\rm reg}(x)=s(1-x^2)/(1+x^2)^2$.
Its omitted jump is handled explicitly at the cusp.
The input to each bulk convolution is extended by zero
below $c$, whereas the output is a vector in full $H$.

The lower endpoint contribution is exactly
\[
 \begin{split}
 \mathcal E_{c,r}v(x)
  ={}&\bigl[a(c)(2\nu r(x-c)-r'(x-c))
                        +a'(c+)r(x-c)\bigr]b(c)\\
     &-a(c)r(x-c)b'(c).
 \end{split}
 \tag{MA-NZ.14}
\]
When $c<0$, splitting the integral at zero adds
$-4\alpha r(x)b(0)$, since $a'(0+)-a'(0-)=-4\alpha$.
At $c=0$ there is only the right-sided endpoint term.
The graph domain in (MA-NZ.4) makes $b,b'$ continuous
across the interior cusp; hence there are no further
terms.  The compact kernel gives zero endpoint terms
at infinity.  These calculations are valid for Sobolev
domain vectors by interval integration by parts.

Apply this identity to $v=R_c^2f$.  Its left side is
$C_rP_c^Zf$.  Thus it is the actual operator equality
\[
 \begin{split}
 C_rP_c^Z={}&\mathcal B_rR_c^2+\mathcal E_{c,r}R_c^2\\
 &-\mathbf1_{c<0}\,4\alpha\,r\otimes
                              [\,f\mapsto(wR_c^2f)(0)\,].
 \end{split}
 \tag{MA-NZ.15}
\]
In particular the last two terms are actual finite-rank maps.

The normalized endpoint functionals
$b(c)/\sqrt{w(c)}$ and $b'(c)/\sqrt{w(c)}$
are bounded uniformly by the second graph norm.
To verify this, apply the ordinary $H^2$ trace inequality
on $[c,c+1]$ and use the uniform comparability of $w$
there with $w(c)$; the three weighted integrals are
exactly the norms of $v,\mathcal L_cv,\mathcal L_c^2v$.
The functional $b(0)$ has the same bound using $[0,1]$
for every $c\le0$.
The uniform graph bound on $R_c^2$ proved after
(MA-NZ.9) therefore makes these normalized composed
functionals uniformly bounded in $c$.

\emph{4. Shrinking the actual test.}
The condition $g(0)=0$ gives, for $r=r_\varepsilon$,
\[
 \begin{gathered}
 \|r''\|_1=O(1),\qquad \int r''=0,\qquad
 \|r'\|_1=O(\varepsilon),\qquad \|r\|_1=O(\varepsilon^2),\\
 \|r'\|_2=O(\varepsilon^{1/2}),\qquad
 \|r\|_2=O(\varepsilon^{3/2}).
 \end{gathered}
 \tag{MA-NZ.16}
\]
For example $g(t)=O(|t|)$ and bounded $g',g''$ on
$|t|\le2\varepsilon$ bound separately the three
terms in $(g\chi(t/\varepsilon))''$.
The cancellation of its integral follows because $r$
is a genuine smooth compact full-line test.
These estimates and the full-line small-shift bound
show $\|\mathcal B_r\|\le C$, uniformly in $c,\varepsilon$.
On a fixed $F_{T,c}$ the principal term satisfies
\[
 C_{r''}v=\int r''(t)(U_tv-v)dt,
 \qquad
 \|\mathcal B_r|_{F_{T,c}}\|
                       \le C_T\sqrt\varepsilon.
 \tag{MA-NZ.17}
\]
Indeed (MA-NZ.12) controls that integral and the other
bulk terms have norms $O(\varepsilon)$ and
$O(\varepsilon^2)$ by (MA-NZ.16).

The range of $R_cE_{T,c}R_c$ is in $F_{T,c}$ and its
trace norm is at most $S_\nu/a_0$.
Using (MA-NZ.9) and (MA-NZ.17) gives
\[
 \sup_{c\le0}\|\mathcal B_rR_c^2\|_1
  \le C a_0^{-1}\sqrt{S_\nu S_\nu(T)}
                   +C_T\sqrt\varepsilon\,S_\nu/a_0.
 \tag{MA-NZ.18}
\]
First let $\varepsilon\downarrow0$ and then $T\to\infty$.
The bulk tends to zero uniformly.

For the moving endpoint, the weighted norms of
$r(\,\cdot-c)$ and $r'(\,\cdot-c)$ are bounded by
$C\sqrt{w(c)}\|r\|_2$ and
$C\sqrt{w(c)}\|r'\|_2$.
Since $a(c)=1/w(c)$ and
$|a'(c+)|\le C/w(c)$, the factor $w(c)$ cancels
against the normalized graph functionals just proved.
Equation (MA-NZ.14) therefore has trace norm
$O(\sqrt\varepsilon)$ after composition with $R_c^2$,
uniformly in $c$.
The interior cusp term has trace norm
$O(\varepsilon^{3/2})$, using the bounded weight near
zero and the uniform functional at zero.
Together with (MA-NZ.15) and (MA-NZ.18), this proves
(MA-NZ.3).
\end{proof}

\paragraph{MA-NZ-002 (the full $g(0)=0$ class has no unaccounted contact term).}
Write $\tau_c(g)=\operatorname{Tr}(C_gP_c^Z)$, the actual
ambient trace of MA-AT-001.
Suppose its limit as $c\to-\infty$ has been proved for
every compact smooth test vanishing on a neighborhood
of zero.  Then it has a unique limit for every
$g\in C_c^\infty(\mathbb R)$ satisfying $g(0)=0$:
\[
 \lim_{c\to-\infty}\tau_c(g)
  =\lim_{\varepsilon\downarrow0}\,\lim_{c\to-\infty}
       \tau_c\bigl(g-r_\varepsilon\bigr).
 \tag{MA-NZ.19}
\]
This extends a proved positive/negative-support formula
without a nonsmooth splitting or a false $C^\infty$
approximation claim.

\begin{proof}
MA-NZ-001 bounds
$|\tau_c(g)-\tau_c(g-r_\varepsilon)|$
by a quantity tending to zero uniformly in $c$.
For every fixed $\varepsilon$ the latter test is
smooth and vanishes near zero, so its limit exists
by hypothesis.  The same uniform estimate for two
values of $\varepsilon$ makes those limits Cauchy.
The elementary three-term estimate then proves
(MA-NZ.19) and its uniqueness.

When the formula away from zero is a finite linear
combination of one-sided Mellin sums over divisors
in fixed vertical strips with $O(\log T)$ local
count, its extension is obtained by the same
one-sided integrals of $g$.
For either half-line, two integrations by parts give
\[
 \left|\int_{\mathbb R_\pm}r_\varepsilon(t)e^{(\sigma+i\gamma)t}dt\right|
   \le C\min\{\varepsilon^2,(1+\gamma^2)^{-1}\},
 \tag{MA-NZ.20}
\]
uniformly on a fixed real-part interval.
At zero the value is zero; the first-derivative
boundary term and the second-derivative integral
are uniformly bounded by (MA-NZ.16).
Summing (MA-NZ.20), split at $|\gamma|=1/\varepsilon$,
gives $O(\varepsilon\log(2/\varepsilon))$.
Thus those actual spectral descriptions pass to
the limit whenever their away-from-zero formula
has separately been established.
\end{proof}

The theorem supplies no negative-time divisor formula
on its own.  It removes the remaining endpoint issue
after such a formula is proved for tests supported
strictly on the negative side.
On the requested space $g(0)=0$, no additional
derivative of a delta distribution can remain.
If one elects to extend a resulting functional to
all tests, two extensions agreeing on this space
differ only by a multiple of evaluation at zero:
subtract $g(0)$ times one fixed test of value one
there to see this directly.  Existence of an
unrenormalized limit for tests with $g(0)\ne0$
is not asserted.

All arithmetic inputs are the actual retained jets,
their exact kernel description and the fixed-function
local count already proved in MA-GW.  The new uniform
argument is the dissipative resolvent factorization,
finite-dimensional triangular estimate and explicit
second-order endpoint calculation above.
It uses no $\mathsf A$, $M$, zero separation,
unproved weighted inner projection or extra
arithmetic positivity principle.
